\subsection{Cyclic proof for reverse-append fusion}
\label{sec:reverse-append-proof}

Figure~\ref{fig:reverse-append-proof} shows the cyclic proof for the
reverse-append identity.  The root goal (bottom) is
$\mathsf{reverse}\;(\mathsf{append}\;\mathit{xs}\;\mathit{ys}) : \mathsf{List}$.
Unfolding $\mathsf{reverse}$ (rule \textsc{fix}) exposes
$\mathsf{revAcc}\;(\mathsf{append}\;\mathit{xs}\;\mathit{ys})\;\mathsf{nil}$.
The LLM-proposed lemma rewrites this to
$\mathsf{revAcc}\;\mathit{ys}\;(\mathsf{reverse}\;\mathit{xs})$ (rule
\textsc{lemma}).  Splitting on $\mathit{ys}$ ($\star$, progress event) yields
two branches.  The nil branch $\mathsf{reverse}\;\mathit{xs}$ is an instance
of the root under the substitution $\mathit{ys}\mapsto\mathsf{nil}$,
witnessed by the \textsc{back} rule.  The cons branch drives further through
structural induction on $\mathit{ys}$.

\begin{figure}[t]
\centering
\footnotesize
\[
\AxiomC{}
\RightLabel{\;\fontsize{6}{7}\selectfont\textsc{nil}}
\UnaryInfC{$\mathsf{reverse}\;\mathit{xs} : \mathsf{List}$}
\qquad
\AxiomC{$\mathit{ys}\mapsto\mathsf{nil}$}
\RightLabel{\;\fontsize{6}{7}\selectfont\textsc{back} \textcolor{red!70}{\textdagger}}
\UnaryInfC{\textcolor{red!60}{$\mathsf{reverse}\;\mathit{xs} : \mathsf{List}$ \textcolor{red!70}{\textdagger}}}
\RightLabel{\;\fontsize{6}{7}\selectfont\textsc{cons}}
\BinaryInfC{$\mathit{ys}:\mathsf{List} \vdash \mathsf{revAcc}\;\mathit{ys}\;(\mathsf{reverse}\;\mathit{xs}) : \mathsf{List}$}
\RightLabel{\;\fontsize{6}{7}\selectfont\textcolor{red!70}{\textsc{split} $\star$}}
\UnaryInfC{$\mathit{ys}:\mathsf{List} \vdash \mathsf{revAcc}\;\mathit{ys}\;(\mathsf{reverse}\;\mathit{xs}) : \mathsf{List}$}
\RightLabel{\;\fontsize{6}{7}\selectfont\textsc{lemma}}
\UnaryInfC{$\mathit{ys}:\mathsf{List} \vdash \mathsf{revAcc}\;(\mathsf{append}\;\mathit{xs}\;\mathit{ys})\;\mathsf{nil} : \mathsf{List}$}
\RightLabel{\;\fontsize{6}{7}\selectfont\textsc{fix}}
\UnaryInfC{$\vdash \mathsf{reverse}\;(\mathsf{append}\;\mathit{xs}\;\mathit{ys}) : \mathsf{List}$ \textcolor{red!70}{\textdagger}}
\DisplayProof
\]

\vspace{-4pt}
{\fontsize{6}{7}\selectfont\textbf{Fig.\ \ref{fig:reverse-append-proof}.} Cyclic proof for reverse-append fusion. Red judgement = companion; $\dagger$ = backlink target; $\star$ = progress event; \textsc{lemma} = LLM-proposed cut.}
\label{fig:reverse-append-proof}
\end{figure}

The structure duplicates the length-map proof exactly: an asynchronous
phase (fix, lemma), a synchronous split (progress $\star$), and a backlink
from the nil branch to the root under $\mathit{ys}\mapsto\mathsf{nil}$.
Without the \textsc{lemma} rule, the step from
$\mathsf{revAcc}\;(\mathsf{append}\;\mathit{xs}\;\mathit{ys})\;\mathsf{nil}$
to $\mathsf{revAcc}\;\mathit{ys}\;(\mathsf{reverse}\;\mathit{xs})$ would
be a stuck generalisation --- the two terms do not anti-unify, and no fold
is possible.  The LLM proposes the lemma as a cut; the sub-prover proves it;
the main proof uses it as a rewrite rule.
